\chapter{The Crucible}

\nopagebreak
\epigraph{"Number theory has traditionally been the purest of pure mathematics. People study number theory for the joy of doing so, not to make money. At least that was largely true until the advent of public key cryptography."}{"Surprising applications of math" (2011)}
\nopagebreak

The recovery from the quiet year didn't look like a return to the past. It looked like something new.

In 2015 and 2016, the volume began to climb back up. But the shape of the posts had changed. The consultant's notebook---scrappy, reactive, brief---was gone. In its place was a practice of sustained explanation.

The years 2013-2017 form the crucible phase: a period when the blog was tested, refined, and transformed. What emerged was fundamentally different from what entered.

\bigskip

\section*{The Shift in Scale}

In the "Notebook" era (2008-2012), the average post was 232 words.
In the "Institution" era (2018-2025), the average post was 448 words.

The blog literally doubled in density.

The shift happened here, in the years following the 2014 dip. It wasn't a sudden jump; it was a gradual deepening.

\begin{longtable}{l p{6.5cm}}
\toprule
\textbf{Year} & \textbf{Average Words} \tabularnewline
\midrule
2014 & 188 \tabularnewline
2015 & 293 \tabularnewline
2016 & 386 \tabularnewline
2017 & 383 \tabularnewline
\bottomrule
\end{longtable}

By 2016, the new standard was set. A typical post was no longer a paragraph; it was an essay. It was substantial enough to contain code, a derivation, and a conclusion.

This change broke the "blogger's trade-off." The conventional wisdom says you can have high frequency or high depth, but not both. If you write every day, you write fluff. If you write essays, you publish once a month.

Cook found a way to do both. By 2018, he was publishing 266 times a year---nearly every weekday---and maintaining that 450-word average. He wasn't choosing between quantity and quality. He was using quantity to \textit{refine} quality.

\bigskip

\section*{The Gamma Function Test}

We can see the deepening by looking at how Cook treats a single topic over time.

Consider the Gamma function, $\Gamma(z)$. It's a generalization of the factorial function to complex numbers. It's notoriously tricky to implement numerically because it grows so fast---$\Gamma(100)$ has 156 digits.

Cook wrote about Gamma in the early years. In 2010, he posted "Diagram of gamma function identities." It was 114 words. It listed a few formulas, provided a link to a reference, and stopped. It was a note to self. \textit{Useful}, but thin.

Compare that to the Gamma posts from the deepening era.

In 2017, he wrote "Approximate inverse of the gamma function" at 481 words. This wasn't a list of formulas. It was a full exploration. He explained the approximation method, provided Python code to implement it, and tested the accuracy across different ranges. He didn't just point to the math; he inhabited it. He translated it.

The pattern repeated across the gamma function posts of this era. Each one deeper than a reference manual, each one building working code, each one explaining not just \textit{what} but \textit{why}.

The 2010 post was a signpost: "This math exists."
The 2017 posts were infrastructure: "Here is how to build with this math."

\bigskip

\section*{The End of the Linklog}

The early internet was built on "linklogs," blogs that primarily pointed to other things. "Check this out." "Read this article."

In 2008, a significant portion of Cook's posts were essentially links.
By 2017, the links were gone.

The blog stopped being a curator of other people's work and became a primary source. This is a critical threshold for any intellectual project. Are you filtering the stream, or are you adding to it?

The deepening meant Cook was no longer relying on external novelty to drive the blog. He was generating the novelty himself, from the inexhaustible mine of mathematics. You don't need to link to the news when you have the Riemann Zeta function. It hasn't changed in 160 years, and it's still more interesting than today's tech headlines.

\bigskip

\section*{The Clinical Trials Extinction}

Focus requires shedding. As the blog deepened into mathematics, it left things behind.

\begin{longtable}{l p{6.5cm}}
\toprule
\textbf{Year} & \textbf{Clinical Trials Posts} \tabularnewline
\midrule
2008 & 24 \tabularnewline
2009 & 6 \tabularnewline
2010 & 6 \tabularnewline
2011 & 2 \tabularnewline
2012 & 6 \tabularnewline
2013 & 1 \tabularnewline
2014 & 1 \tabularnewline
2015 & 5 \tabularnewline
2016 & 3 \tabularnewline
2018 & 1 \tabularnewline
\bottomrule
\end{longtable}

The pattern is extinction. Two dozen posts in 2008, then a steady decline to single digits, then effectively zero in recent years.

Clinical trials was Cook's day job at MD Anderson, where he spent thirteen years---first as a biostatistician designing trials, later in software development management. By the end, he was managing developers rather than doing mathematics. When he left in 2013 to found Singular Value Consulting, the topic began to fade from the blog. By 2018, it was down to a single post.

He has 5,881 Google Scholar citations, mostly from that clinical trials work. His paper on "dose-finding based on efficacy-toxicity trade-offs" remains influential in biostatistics. He could have built an audience around clinical trial methodology.

He didn't. The topic faded because it was professional obligation, not personal fascination. The business posts tell a similar story: 126 posts in the early years, many reflecting the frustrations of corporate software management, dwindling to 27 by the recent era. Going independent meant he could finally write about what he wanted to think about.

\bigskip

\section*{The PowerShell Fossil Record}

PowerShell tells a similar story.

Cook wrote about PowerShell consistently from 2008-2012, producing 47 posts on scripting techniques, language features, and Windows automation. This represents serious investment in mastering the tool.

After 2012, PowerShell posts decline. The last PowerShell post appears in July 2014. The tool had been replaced.

The replacement was Python. The first SciPy-tagged post appeared in February 2009. By 2025, Python had accumulated 207 tagged posts. The transition was generational: Windows automation gave way to scientific computing.

\bookfigure{ch05_languages_stacked.png}{The language transition: PowerShell fades as Python rises. By 2018, Python dominates the blog's code examples.}

\bigskip

\section*{What Rose as Others Fell}

The shedding wasn't just loss. It created space.

\begin{longtable}{p{2.25cm}p{2.25cm}p{2.25cm}p{2.25cm}}
\toprule
\textbf{Category} & \textbf{2008-2012} & \textbf{2018-2025} & \textbf{Change} \tabularnewline
\midrule
Math & 357 (21\%) & 1,611 (64\%) & +43 pts \tabularnewline
Statistics & 205 (12\%) & 120 (5\%) & -7 pts \tabularnewline
Computing & 145 (8\%) & 314 (12\%) & +4 pts \tabularnewline
Business & 126 (7\%) & 27 (1\%) & -6 pts \tabularnewline
Clinical trials & 44 (3\%) & 1 (0\%) & -3 pts \tabularnewline
\bottomrule
\end{longtable}

Math went from about one-fifth to nearly two-thirds of output. The narrowing was dramatic.

This isn't accidental. Every hour spent on business advice was an hour not spent on number theory. Every PowerShell post was a Python post not written. The later blog's mathematical depth required abandoning breadth.

\bigskip

\section*{Quality Through Repetition}

How do you double your output while doubling your depth?

The answer lies in the subject matter. Mathematics rewards repetition.

If you write about politics, yesterday's post is trash today. If you write about software frameworks, your tutorial from 2015 is broken by 2017.

But if you write about probability distributions, the work compounds. A post about the Beta distribution written in 2015 helps you write a post about the Dirichlet distribution in 2016. The code you wrote for numerical integration in 2014 becomes the library you import in 2017.

The deepening was possible because the blog found its true subject. Cook moved away from the ephemeral (tech news, business advice) and anchored himself in the permanent.

He wasn't just writing more. He was building a cathedral, brick by brick. The foundation laid in the quiet year supported the weight of the structure to come.

\bigskip

\section*{What Persisted}

Some things didn't change:

\textbf{Pi} (153 posts) spans all seventeen years. The first pi post appeared in 2008; the most recent in 2025.

\textbf{Primes} (103 posts) similarly persist. The mathematical core remained stable while everything around it shifted.

\textbf{The Fourier transform} (45 tagged posts) appears throughout. Signal processing was never a professional obligation; it was genuine interest.

\textbf{Special functions} (140 posts) never fade. The gamma function, Bessel functions, hypergeometric functions---these define the blog's mathematical identity.

The pattern: professional topics faded. Personal fascinations persisted. The blog became a map of what Cook actually wanted to think about, freed from what he was paid to think about.

\bigskip

\section*{So What?}

\textit{The crucible phase refined the blog into its final form. Posts doubled in depth while professional obligations burned away. Clinical trials went to zero. PowerShell became a fossil. Business advice faded. What emerged was a blog about mathematics---not because the other topics failed, but because the crucible revealed what Cook actually wanted to think about when nobody was paying him to think. The transition from 21\% mathematics to 64\% mathematics wasn't narrowing. It was clarification.}

